% ============================================================
\section{Defining CM and DM: From Words to Mathematics}
% ============================================================

In Section~0 we described DM and CM in plain language.  Now we turn
those descriptions into precise mathematics — definitions that can be
applied to any number of lines and packaged into a matrix
transformation that the rest of the paper uses as its engine.

% ------------------------------------------------------------
\subsection{The General \texorpdfstring{$N$}{N}-Line System}
% ------------------------------------------------------------

\figentry{fig:P1Fig1}{GDCMDMCEAP_Fig1_page01}
{General $N$-line model used to define CM and DM quantities.
 Each line carries current $i_n$ through a series $R$--$L$ impedance.
 Shunt capacitors connect each line to ground: $C$ for lines 1 through
 $N-1$, and $kC$ for line $N$ (the asymmetric line).
 The floating reference $P$ is distinct from ground;
 $v_{nP}$ is the voltage from line $n$ to $P$,
 and $v_{Pg}$ is the voltage from $P$ to ground.}{0.72}

\figref{fig:P1Fig1} is the general model the paper uses throughout.
Walk through it from left to right:

\begin{circuitbox}[Circuit Walk — Fig.~1]
\begin{itemize}
    \item \textbf{$N$ parallel lines:} Each horizontal wire is one
          ``line'' — it could be a phase conductor, a terminal of a
          transistor chip, or a filter inductor branch.  Line $n$
          carries current $i_n$ from left to right ($i_n'$ exits on
          the far right after passing through the series $R$--$L$).

    \item \textbf{Series $R$--$L$ impedance:} Each line has a
          resistance $R$ and inductance $L$ in series.  This models
          cable resistance, bond-wire inductance, or busbar
          impedance — whatever is in series with that conductor.

    \item \textbf{Shunt capacitors to ground:} On the right end,
          each line has a capacitor to ground.  Lines 1 through $N-1$
          each have capacitance $C$.  Line $N$ has capacitance $kC$,
          where $k \neq 1$ makes this line asymmetric with respect to
          the others.  In the MCPM application ($k < 1$), this is the
          source terminal whose capacitance to the baseplate is smaller
          than the rest.

    \item \textbf{Floating reference $P$:} On the left, a point $P$
          is defined with voltages $v_{nP}$ measured from each line $n$
          to $P$.  Crucially, $P$ is \emph{not} ground — it floats.
          Its voltage with respect to ground is $v_{Pg}$.
\end{itemize}
\end{circuitbox}

% ------------------------------------------------------------
\subsection{Why a Floating Reference Instead of Ground?}
% ------------------------------------------------------------

In classical two-wire analysis, voltages are measured to ground and
DM is simply the full line-to-line voltage ($v_1 - v_2$).  The paper
adopts a more general approach: voltages are measured to an
\emph{arbitrary floating reference} $P$.  This may seem like unnecessary abstraction,
but it has a powerful consequence that becomes clear in the
decomposition:

\begin{physbox}[The Key Result of Using a Floating Reference]
After decomposition, the ground voltage $v_{Pg}$ appears
\textbf{only} in the CM equation — it vanishes from every DM
equation.  This means the DM quantities are completely
\emph{independent of ground}.  Only the CM equation ``sees'' the
ground connection.

Practically, $P$ can be chosen to be any convenient internal circuit
node (such as the phase terminal~$A$ of the half-bridge in \S4).
That choice replaces the abstract $v_{Pg}$ with a known circuit
voltage $v_{Ag}$ — the half-bridge switching waveform — turning
the CM equation's right-hand side into a concrete, measurable
driving source.
\end{physbox}

% ------------------------------------------------------------
\subsection{CM Current and Voltage: The Definitions}
% ------------------------------------------------------------

For the $N$-line system of \figref{fig:P1Fig1}, the CM quantities are:

\begin{eqnbox}[CM Definitions — Equations (1) and (2)]
\begin{align}
    \icm  &\triangleq \sum_{n=1}^{N} i_n
           \label{eq:icm} \\[4pt]
    \vcm  &\triangleq \frac{1}{N} \sum_{n=1}^{N} v_{nP}
           \label{eq:vcm}
\end{align}
\textbf{CM current} is the sum of all line currents.
\textbf{CM voltage} is the average of all line-to-$P$ voltages.

\medskip
\textbf{Why a sum for current but an average for voltage?}

\textit{Current:} By KCL at the ground node, the current that
actually flows to ground equals $i_1 + i_2 + \cdots + i_N$ — the
total imbalance.  A clamp-meter placed around all $N$ wires reads
exactly this sum (DM currents cancel; CM currents add).  Dividing
by $N$ would give a per-line average that no single instrument
measures and that has no direct circuit meaning.

\textit{Voltage:} The CM voltage should capture the
``common level'' shared by all lines — the average is the natural
definition.  The sum would grow with $N$ and lose physical meaning.

\textit{The deeper reason — impedance consistency:}  For $N$
symmetric lines each with resistance $R$, CM current divides
equally across $N$ parallel branches, so the correct CM resistance
is $R/N$.  With $\vcm = \frac{1}{N}\sum v_{nP}$ and
$\icm = \sum i_n$, Ohm's law gives
$Z_{\mathrm{CM}} = \vcm/\icm = R/N$ \checkmark.
If instead $\icm$ were the average, $Z_{\mathrm{CM}}$ would
come out as $R$ — as if $N$ parallel resistors were a single wire.
The asymmetric definitions are precisely what keeps the transformed
circuit elements physically correct.
\end{eqnbox}

These match the intuition from \S0: CM current is what you measure
with a clamp around all wires together, and CM voltage is the
``average'' floating potential of all lines relative to the chosen
reference.

\begin{warningbox}[CM Current Is Signed]
In a perfectly balanced system all line currents sum to zero:
$\icm = 0$.  Any non-zero $\icm$ is entirely due to currents that
have ``escaped'' through the shunt capacitors to ground — exactly
the conducted EMI we are trying to model.
\end{warningbox}

% ------------------------------------------------------------
\subsection{DM Current and Voltage: The Definitions}
% ------------------------------------------------------------

The paper defines DM quantities for every \emph{adjacent pair} of
lines.  For lines $m$ and $n$:

\begin{eqnbox}[DM Definitions — Equations (3) and (4)]
\begin{align}
    i_{mn} &\triangleq \frac{1}{\Ni}(i_m - i_n)
            \label{eq:idm} \\[4pt]
    v_{mn} &\triangleq v_{mP} - v_{nP}
            \label{eq:vdm}
\end{align}
where $\Ni$ is a normalisation factor.
\end{eqnbox}

For $N$ lines there are $N-1$ adjacent DM pairs
(1--2, 2--3, \ldots, $(N-1)$--$N$) plus 1 CM quantity, giving $N$
quantities total — the same count as the original $N$ line currents.
The set is therefore a complete, invertible re-description of the
original variables.

\textbf{Why does $\Ni$ exist, and why is it 2?}

The normalisation factor $\Ni$ is a free parameter — you could set it
to 1, to $\sqrt{N}$, or to any non-zero value and the mathematics
would remain valid.  The question is: which choice produces the most
\emph{physically readable} decomposed circuit?

Consider two symmetric lines, each with resistance $R$, carrying
a pure DM current of 3\,A (3\,A out on line~1, 3\,A back on line~2).
Trace the DM current path: it travels through $R$ on line~1, then
back through $R$ on line~2.  The total loop resistance is $2R$.
You would draw this DM loop as a single resistor of value $2R$.

Now compute $i_{12}$ with each candidate $\Ni$:

\begin{center}
\renewcommand{\arraystretch}{1.4}
\begin{tabular}{cccc}
\toprule
$\Ni$ & $i_{12} = \frac{i_1 - i_2}{\Ni}$ & DM element & Physical match? \\
\midrule
$1$      & $\frac{3-(-3)}{1} = 6\,\text{A}$  & $1\cdot R = R$       & $6\,\text{A}$ through $R$? No — actual wire carries 3\,A \\
$2$      & $\frac{3-(-3)}{2} = 3\,\text{A}$  & $2\cdot R = 2R$      & $3\,\text{A}$ through $2R$ = series loop \checkmark \\
$\sqrt{2}$ & $\frac{6}{\sqrt{2}} \approx 4.2\,\text{A}$ & $\sqrt{2}R$ & $4.2\,\text{A}$ through $\sqrt{2}R$? No clear meaning \\
\bottomrule
\end{tabular}
\end{center}

\begin{physbox}[Why $\Ni = 2$ is the right choice]
With $\Ni = 2$ the DM current $i_{12} = (i_1 - i_2)/2$ equals the
actual physical current flowing in each wire during pure DM
($i_{12} = 3\,\text{A}$ = the actual 3\,A in line~1).  As a result,
the DM impedance that appears in the decomposed equations is exactly
the \emph{series combination} of the two per-line impedances — the
same value you would assign if you physically drew the DM current
loop and added up the resistances it passes through.

With $\Ni = 1$ the ``DM current'' is 6\,A but each wire only carries
3\,A — the number has no direct physical counterpart.
With $\Ni = \sqrt{2}$ the DM impedance $\sqrt{2}\,R$ has no simple
circuit interpretation at all.

Only $\Ni = 2$ keeps the DM current equal to the per-wire current
\emph{and} the DM impedance equal to the familiar series combination.
This is the sole reason for the choice.
\end{physbox}

% ------------------------------------------------------------
\subsection{Building the Transformation Matrices: \texorpdfstring{$2{\times}2$}{2x2},
            \texorpdfstring{$3{\times}3$}{3x3}, then General}
% ------------------------------------------------------------

A matrix is simply a compact way of writing several equations at once.
Rather than presenting the general $N \times N$ matrix immediately,
we build it up one size at a time so the structure is obvious before
the notation becomes dense.

% ---- Step 1: N = 2 ----
\subsubsection*{Step 1 — Two lines ($N = 2$)}

With 2 lines we need to encode 2 quantities: one DM (pair 1--2)
and one CM.  Before writing the matrix, look at the circuit
whose variables it encodes.

\begin{figure}[H]
\centering
\begin{circuitikz}[american, scale=0.82, transform shape]

  %--- Line 1 (y = 2) ---
  \draw (0,2) node[ocirc]{}
        to[R, l=$R$]  (2.5,2)
        to[L, l=$L$]  (5.0,2)
        -- (8.5,2);
  \fill[black] (6.0,2) circle (2.2pt);               % branch junction for C_1
  \draw[->, thick] (0.25,2.38) -- (1.1,2.38)
        node[midway, above, font=\small]{$i_1$};
  \draw[->, thick] (7.0,2.38) -- (8.0,2.38)
        node[midway, above, font=\small]{$i_1'$};

  %--- Line 2 (y = 0) ---
  \draw (0,0) node[ocirc]{}
        to[R, l_=$R$] (2.5,0)
        to[L, l_=$L$] (5.0,0)
        -- (8.5,0);
  \fill[black] (6.8,0) circle (2.2pt);               % branch junction for C_2
  \draw[->, thick] (0.25,-0.38) -- (1.1,-0.38)
        node[midway, below, font=\small]{$i_2$};
  \draw[->, thick] (7.5,-0.38) -- (8.3,-0.38)
        node[midway, below, font=\small]{$i_2'$};

  %--- Shunt capacitors — common ground bus at y = -2.0 ---
  % C_1 at x=6.0: symbol sits between y=2 and y=0.7 (clearly above line 2);
  % plain wire from y=0.7 to y=-2.0 crosses line 2 without connection (no junction dot)
  \draw (6.0,2) to[C, l=$C$] (6.0,0.7) -- (6.0,-2.0);
  % C_2 at x=6.8: symbol sits between y=0 and y=-2.0 (no crossings)
  \draw (6.8,0) to[C, l=$C$] (6.8,-2.0);
  % Common ground bus
  \draw (5.4,-2.0) -- (7.5,-2.0);
  \draw (6.4,-2.0) node[ground]{};

  %--- Floating reference P — NO physical wire to ground ---
  \coordinate (Pnd) at (-1.5, 1.0);
  \fill[black] (Pnd) circle (2.5pt);
  \node[left, font=\small\bfseries] at (Pnd) {$P$};

  %--- v_{1P}: dashed voltage label from line 1 terminal to P ---
  \draw[dashed, thin] (0,2) -- (Pnd)
        node[pos=0.62, right, font=\scriptsize]{$v_{1P}$};

  %--- v_{2P}: dashed voltage label from line 2 terminal to P ---
  \draw[dashed, thin] (0,0) -- (Pnd)
        node[pos=0.38, right, font=\scriptsize]{$v_{2P}$};

  %--- v_{Pg}: dashed — P is floating; dashed = voltage annotation, not a wire ---
  \draw[dashed, thin] (Pnd) -- (-1.5,-2.0) node[ground]{};
  \node[right, font=\scriptsize] at (-1.5,-0.3) {$v_{Pg}$};

  %--- Line labels ---
  \node[above, font=\footnotesize, color=gray!70!black] at (0,2) {Line 1};
  \node[below, font=\footnotesize, color=gray!70!black] at (0,0) {Line 2};

\end{circuitikz}
\caption{Circuit for $N = 2$ lines.
  Column vectors: $\ivec = [i_1,\;i_2]^T$ and $\vvec = [v_{1P},\;v_{2P}]^T$.
  \textbf{$P$ is a floating reference} — no physical wire connects it to ground.
  Dashed lines (including the one labeled $v_{Pg}$) are voltage annotations, not conductors.
  Both shunt capacitors share a single common ground bus (solid horizontal bar at bottom).
  Primed currents $i_n'$ exit the right side after the shunt-capacitor junction.
  A wire crossing without a filled dot indicates no electrical connection.}
\label{fig:ckt_N2}
\end{figure}

The matrix $\mathbf{T}^i_2$ reads the two currents $[i_1,\,i_2]^T$
and produces the DM current $i_{12}$ and CM current $\icm$.
The matrix $\mathbf{T}^v_2$ reads the two voltages $[v_{1P},\,v_{2P}]^T$
and produces the DM voltage $v_{12}$ and CM voltage $\vcm$:

\[
\mathbf{T}^{i}_{2} =
\begin{bmatrix}
\tfrac{1}{2} & -\tfrac{1}{2} \\[4pt]
1 & 1
\end{bmatrix},
\qquad
\mathbf{T}^{v}_{2} =
\begin{bmatrix}
1 & -1 \\[4pt]
\tfrac{1}{2} & \tfrac{1}{2}
\end{bmatrix}.
\]

Multiply out to verify:
\[
\mathbf{T}^{i}_{2}
\begin{bmatrix} i_1 \\ i_2 \end{bmatrix}
=
\begin{bmatrix} \tfrac{i_1 - i_2}{2} \\ i_1 + i_2 \end{bmatrix}
=
\begin{bmatrix} i_{12} \\ \icm \end{bmatrix},
\qquad
\mathbf{T}^{v}_{2}
\begin{bmatrix} v_{1P} \\ v_{2P} \end{bmatrix}
=
\begin{bmatrix} v_{1P} - v_{2P} \\ \tfrac{v_{1P}+v_{2P}}{2} \end{bmatrix}
=
\begin{bmatrix} v_{12} \\ \vcm \end{bmatrix}.
\]

Row~1 takes the scaled difference of the two currents
(the DM definition).  Row~2 sums them (the CM definition).
The two rows together account for all the information in $i_1$ and
$i_2$ — nothing is lost.

% ---- Step 2: N = 3 ----
\subsubsection*{Step 2 — Three lines ($N = 3$)}

With 3 lines we need to encode 3 quantities: DM pair 1--2,
DM pair 2--3, and one CM.  The circuit gains a third line;
note that line~3 carries capacitance $kC$ ($k \neq 1$),
making it the asymmetric line.

\begin{figure}[H]
\centering
\begin{circuitikz}[american, scale=0.82, transform shape]

  %--- Line 1 (y = 4) ---
  \draw (0,4) node[ocirc]{}
        to[R, l=$R$]  (2.5,4)
        to[L, l=$L$]  (5.0,4)
        -- (8.5,4);
  \fill[black] (6.0,4) circle (2.2pt);               % branch junction for C_1
  \draw[->, thick] (0.25,4.38) -- (1.1,4.38)
        node[midway, above, font=\small]{$i_1$};
  \draw[->, thick] (7.0,4.38) -- (8.0,4.38)
        node[midway, above, font=\small]{$i_1'$};

  %--- Line 2 (y = 2) ---
  \draw (0,2) node[ocirc]{}
        to[R, l=$R$]  (2.5,2)
        to[L, l=$L$]  (5.0,2)
        -- (8.5,2);
  \fill[black] (6.8,2) circle (2.2pt);               % branch junction for C_2
  \draw[->, thick] (0.25,1.62) -- (1.1,1.62)
        node[midway, below, font=\small]{$i_2$};
  \draw[->, thick] (7.5,1.62) -- (8.3,1.62)
        node[midway, below, font=\small]{$i_2'$};

  %--- Line 3 (y = 0) — asymmetric capacitance kC ---
  \draw (0,0) node[ocirc]{}
        to[R, l_=$R$] (2.5,0)
        to[L, l_=$L$] (5.0,0)
        -- (8.5,0);
  \fill[black] (7.6,0) circle (2.2pt);               % branch junction for C_3
  \draw[->, thick] (0.25,-0.38) -- (1.1,-0.38)
        node[midway, below, font=\small]{$i_3$};
  \draw[->, thick] (8.0,-0.38) -- (8.5,-0.38)
        node[midway, below, font=\small]{$i_3'$};

  %--- Shunt capacitors — common ground bus at y = -2.5 ---
  % C_1 at x=6.0: symbol between y=4 and y=2.8 (above line 2);
  % plain wire from y=2.8 down to bus crosses lines 2 and 3 (no junction dot = no connection)
  \draw (6.0,4) to[C, l=$C$]   (6.0,2.8) -- (6.0,-2.5);
  % C_2 at x=6.8: symbol between y=2 and y=0.8 (above line 3);
  % plain wire from y=0.8 down to bus crosses line 3 (no junction dot = no connection)
  \draw (6.8,2) to[C, l=$C$]   (6.8,0.8) -- (6.8,-2.5);
  % C_3 at x=7.6: symbol between y=0 and y=-2.5 (no crossings)
  \draw (7.6,0) to[C, l_=$kC$] (7.6,-2.5);
  % Common ground bus
  \draw (5.3,-2.5) -- (8.2,-2.5);
  \draw (6.8,-2.5) node[ground]{};

  %--- Floating reference P — NO physical wire to ground ---
  \coordinate (Pnd) at (-1.5, 2.0);
  \fill[black] (Pnd) circle (2.5pt);
  \node[left, font=\small\bfseries] at (Pnd) {$P$};

  %--- v_{1P}, v_{2P}, v_{3P}: dashed voltage labels ---
  \draw[dashed, thin] (0,4) -- (Pnd)
        node[pos=0.65, right, font=\scriptsize]{$v_{1P}$};
  \draw[dashed, thin] (0,2) -- (Pnd)
        node[midway, above, font=\scriptsize]{$v_{2P}$};
  \draw[dashed, thin] (0,0) -- (Pnd)
        node[pos=0.35, right, font=\scriptsize]{$v_{3P}$};

  %--- v_{Pg}: dashed — P is floating; dashed = voltage annotation, not a wire ---
  \draw[dashed, thin] (Pnd) -- (-1.5,-2.5) node[ground]{};
  \node[right, font=\scriptsize] at (-1.5, 0.6) {$v_{Pg}$};

  %--- Line labels ---
  \node[above, font=\footnotesize, color=gray!70!black] at (0,4) {Line 1};
  \node[above, font=\footnotesize, color=gray!70!black] at (0,2) {Line 2};
  \node[below, font=\footnotesize, color=gray!70!black] at (0,0) {Line 3\;($kC$, asymmetric)};

\end{circuitikz}
\caption{Circuit for $N = 3$ lines.
  Line~3 is asymmetric: shunt capacitance $kC$ with $k < 1$;
  lines~1 and 2 each have capacitance~$C$.
  Input vectors: $\ivec = [i_1,\;i_2,\;i_3]^T$ and $\vvec = [v_{1P},\;v_{2P},\;v_{3P}]^T$.
  \textbf{$P$ is floating} (dashed lines are voltage annotations, not conductors).
  All three caps share a single common ground bus.
  Wire crossings without a filled dot indicate no connection.}
\label{fig:ckt_N3}
\end{figure}

Adding a third line means the matrices grow to $3 \times 3$:

\[
\mathbf{T}^{i}_{3} =
\begin{bmatrix}
\tfrac{1}{2} & -\tfrac{1}{2} & 0 \\[4pt]
0 & \tfrac{1}{2} & -\tfrac{1}{2} \\[4pt]
1 & 1 & 1
\end{bmatrix},
\qquad
\mathbf{T}^{v}_{3} =
\begin{bmatrix}
1 & -1 & 0 \\[4pt]
0 & 1 & -1 \\[4pt]
\tfrac{1}{3} & \tfrac{1}{3} & \tfrac{1}{3}
\end{bmatrix}.
\]

Multiply out:
\[
\mathbf{T}^{i}_{3}
\begin{bmatrix} i_1 \\ i_2 \\ i_3 \end{bmatrix}
=
\begin{bmatrix}
\tfrac{i_1 - i_2}{2} \\[3pt]
\tfrac{i_2 - i_3}{2} \\[3pt]
i_1 + i_2 + i_3
\end{bmatrix}
=
\begin{bmatrix} i_{12} \\ i_{23} \\ \icm \end{bmatrix},
\qquad
\mathbf{T}^{v}_{3}
\begin{bmatrix} v_{1P} \\ v_{2P} \\ v_{3P} \end{bmatrix}
=
\begin{bmatrix}
v_{1P} - v_{2P} \\[3pt]
v_{2P} - v_{3P} \\[3pt]
\tfrac{v_{1P}+v_{2P}+v_{3P}}{3}
\end{bmatrix}
=
\begin{bmatrix} v_{12} \\ v_{23} \\ \vcm \end{bmatrix}.
\]

\begin{physbox}[Spotting the Pattern Across $N = 2$ and $N = 3$]
Compare the two matrices side by side and notice:
\begin{itemize}
    \item \textbf{DM rows (all rows except the last):} Each DM row
          has $+\tfrac{1}{2}$ on one column and $-\tfrac{1}{2}$ on
          the next column (for $\TiN$), or $+1$ and $-1$
          (for $\TvN$), and zeros everywhere else.  Row $k$ always
          captures the \emph{adjacent} pair $(k,\, k+1)$.

    \item \textbf{CM row (always the last row):} For $\TiN$ it is
          all ones — sum all currents.  For $\TvN$ it is all
          $\tfrac{1}{N}$ — average all voltages.

    \item \textbf{Growing the matrix for larger $N$:} Adding a
          fourth line just inserts one more DM row (for pair 3--4)
          above the CM row, and shifts the CM-row weights from
          $\tfrac{1}{3}$ to $\tfrac{1}{4}$.  The structure never
          changes — only the size.
\end{itemize}
\end{physbox}

% ---- Step 3: General N ----
\subsubsection*{Step 3 — General $N$ lines (Equations (5)--(8) in the paper)}

Now that the pattern is clear, the general form writes itself.
Collect all line currents into
$\ivec = [i_1,\,\ldots,\,i_N]^T$ and all decomposed quantities into
$\iDCM = [i_{12},\,i_{23},\,\ldots,\,i_{(N-1)N},\,\icm]^T$,
and likewise for voltages.  Then:

\begin{eqnbox}[Transformation Equations — Equations (5) and (6)]
\begin{align}
    \iDCM &= \TiN\, \ivec  \label{eq:Ti} \\[4pt]
    \vDCM &= \TvN\, \vvec  \label{eq:Tv}
\end{align}
\end{eqnbox}

where the $N \times N$ matrices are (Equations (7) and (8)):

\begin{equation}
\label{eq:TiN}
\TiN =
\begin{bmatrix}
\frac{1}{\Ni} & -\frac{1}{\Ni} & 0 & \cdots & 0 & 0 \\[4pt]
0 & \frac{1}{\Ni} & -\frac{1}{\Ni} & \cdots & 0 & 0 \\
\vdots & & \ddots & \ddots & & \vdots \\
0 & 0 & \cdots & 0 & \frac{1}{\Ni} & -\frac{1}{\Ni} \\[4pt]
1 & 1 & \cdots & 1 & 1 & 1
\end{bmatrix}
\end{equation}

\begin{equation}
\label{eq:TvN}
\TvN =
\begin{bmatrix}
1 & -1 & 0 & \cdots & 0 & 0 \\[4pt]
0 & 1 & -1 & \cdots & 0 & 0 \\
\vdots & & \ddots & \ddots & & \vdots \\
0 & 0 & \cdots & 0 & 1 & -1 \\[4pt]
\frac{1}{N} & \frac{1}{N} & \cdots & \frac{1}{N} & \frac{1}{N} & \frac{1}{N}
\end{bmatrix}
\end{equation}

Both matrices are square and invertible — the original $N$ line
variables can always be recovered from the decomposed set.

% ------------------------------------------------------------
\subsection{What the Inverse Gives You}
% ------------------------------------------------------------

The matrices $\TiN$ and $\TvN$ are invertible, so the mapping works
in both directions: given the DM and CM quantities you can always
recover the original per-line variables.  This section derives the
recovery formulas for $N = 2$ and works through a numerical example
for both currents and voltages.

% ---- current inverse ----
\subsubsection*{Deriving the Current Inverse ($N = 2$)}

Start from the two forward equations:
\[
    i_{12} = \frac{i_1 - i_2}{2}
    \qquad\text{and}\qquad
    \icm   = i_1 + i_2.
\]
Rewrite the first as $i_1 - i_2 = 2\,i_{12}$.
Add this to the second equation ($i_1 + i_2 = \icm$):
\[
    2\,i_1 = \icm + 2\,i_{12}
    \;\;\Longrightarrow\;\;
    \boxed{i_1 = \frac{\icm}{2} + i_{12}.}
\]
Subtract the rearranged first from the second:
\[
    2\,i_2 = \icm - 2\,i_{12}
    \;\;\Longrightarrow\;\;
    \boxed{i_2 = \frac{\icm}{2} - i_{12}.}
\]

\begin{eqnbox}[Recovering Line Currents from DM and CM ($N = 2$)]
\begin{align}
    i_1 &= \frac{\icm}{2} + i_{12} \label{eq:inv_i1} \\[4pt]
    i_2 &= \frac{\icm}{2} - i_{12} \label{eq:inv_i2}
\end{align}
Each wire carries the \textbf{same CM share} ($\icm/2$: the total CM
current distributed equally between the two lines) plus or minus the
\textbf{DM current} $i_{12}$ (equal-and-opposite push-pull component).
\end{eqnbox}

% ---- voltage inverse ----
\subsubsection*{Deriving the Voltage Inverse ($N = 2$)}

The forward voltage equations are:
\[
    v_{12} = v_{1P} - v_{2P}
    \qquad\text{and}\qquad
    \vcm   = \frac{v_{1P} + v_{2P}}{2}.
\]
Rewrite the second as $v_{1P} + v_{2P} = 2\,\vcm$.
Add this to the first ($v_{1P} - v_{2P} = v_{12}$):
\[
    2\,v_{1P} = 2\,\vcm + v_{12}
    \;\;\Longrightarrow\;\;
    \boxed{v_{1P} = \vcm + \frac{v_{12}}{2}.}
\]
Subtract the first from the second rearrangement:
\[
    2\,v_{2P} = 2\,\vcm - v_{12}
    \;\;\Longrightarrow\;\;
    \boxed{v_{2P} = \vcm - \frac{v_{12}}{2}.}
\]

\begin{eqnbox}[Recovering Line Voltages from DM and CM ($N = 2$)]
\begin{align}
    v_{1P} &= \vcm + \frac{v_{12}}{2} \label{eq:inv_v1P} \\[4pt]
    v_{2P} &= \vcm - \frac{v_{12}}{2} \label{eq:inv_v2P}
\end{align}
Each line-to-$P$ voltage is the shared CM voltage (common to both
wires) plus or minus half the DM voltage.
\end{eqnbox}

% ---- worked example: currents ----
\subsubsection*{Worked Example — Currents}

\textbf{Scenario.}  A two-wire cable carries 8\,A of DC load current
to a resistive load.  Parasitic module capacitances leak a total of
0.4\,A of common-mode current to ground.  Find the ammeter reading
on each wire and verify.

\textbf{Given:} $i_{12} = 8\,\text{A}$,\quad $\icm = 0.4\,\text{A}$.

\textbf{Step 1 — apply the inverse:}
\[
    i_1 = \frac{\icm}{2} + i_{12}
        = \frac{0.4}{2} + 8
        = 0.2 + 8
        = 8.2\,\text{A}
\]
\[
    i_2 = \frac{\icm}{2} - i_{12}
        = 0.2 - 8
        = -7.8\,\text{A}
\]

\textbf{Step 2 — verify by re-applying the forward transform:}
\[
    i_{12} = \frac{i_1 - i_2}{2}
           = \frac{8.2 - (-7.8)}{2}
           = \frac{16}{2}
           = 8\,\text{A} \checkmark
\]
\[
    \icm = i_1 + i_2 = 8.2 + (-7.8) = 0.4\,\text{A} \checkmark
\]

\begin{physbox}[Reading the Numbers]
Line~1 carries 8.2\,A forward and line~2 carries 7.8\,A return.
The 0.4\,A imbalance ($8.2 - 7.8 = 0.4$\,A) is $\icm$: it left
line~1 and did not return on line~2 — it leaked to ground through
the parasitic capacitance.  This is exactly the conducted-EMI
current that the LISN would measure.  In a perfectly symmetric
system ($\icm = 0$) both wires would read exactly $\pm\,8$\,A.
\end{physbox}

% ---- worked example: voltages ----
\subsubsection*{Worked Example — Voltages}

\textbf{Scenario.}  The same cable has a 600\,V line-to-line voltage
($v_{12} = 600\,\text{V}$) and a common-mode voltage
$\vcm = 300\,\text{V}$ (the average of both lines sits 300\,V above
the floating reference~$P$).  Find $v_{1P}$ and $v_{2P}$.

\textbf{Given:} $v_{12} = 600\,\text{V}$,\quad $\vcm = 300\,\text{V}$.

\textbf{Step 1 — apply the inverse:}
\[
    v_{1P} = \vcm + \frac{v_{12}}{2}
           = 300 + \frac{600}{2}
           = 300 + 300
           = 600\,\text{V}
\]
\[
    v_{2P} = \vcm - \frac{v_{12}}{2}
           = 300 - 300
           = 0\,\text{V}
\]

\textbf{Step 2 — verify:}
\[
    v_{12} = v_{1P} - v_{2P} = 600 - 0 = 600\,\text{V} \checkmark
\]
\[
    \vcm = \frac{v_{1P} + v_{2P}}{2}
         = \frac{600 + 0}{2}
         = 300\,\text{V} \checkmark
\]

\begin{physbox}[Why $v_{2P} = 0$\,V?]
Line~2 coincides with the floating reference~$P$ in this example
($v_{2P} = 0$) — not because of any special circuit property, but
because $\vcm = v_{12}/2$ for the chosen numbers.  A different CM
voltage would place $P$ at a different level between the two lines.
Notice that neither calculation required knowing $v_{Pg}$ (the
voltage from $P$ to ground): once $P$ is chosen, all DM and CM
quantities are determined independently of ground.
\end{physbox}

\begin{warningbox}[The Factor of $N$ in the CM Share]
Because $\icm = i_1 + i_2$ is the \emph{total} CM current
(the sum, per Eq.~\ref{eq:icm}), the per-line CM contribution is
$\icm/N = \icm/2$ for $N = 2$.  A common mistake is to write
$i_1 = \icm + i_{12}$, forgetting the $1/2$.  Check: with that
wrong formula $i_1 + i_2 = 2\icm \neq \icm$ — the CM current
double-counts.  The correct rule is:
\emph{the total CM current is shared equally among all $N$ lines.}
\end{warningbox}

\begin{statebox}[Summary of \S1]
\begin{itemize}
    \item \textbf{CM:} sum of line currents; average of line-to-$P$
          voltages. Appears in conducted EMI. Connects to ground
          through $v_{Pg}$.
    \item \textbf{DM ($N-1$ pairs):} scaled difference of adjacent
          line currents; line-to-line voltage differences.
          Independent of ground.
    \item \textbf{Transformation matrices $\TiN$, $\TvN$:}
          $N \times N$, invertible, pack all CM/DM definitions into
          two matrix equations.
    \item \textbf{Normalisation $\Ni = 2$:} gives DM circuit elements
          equal to series combinations of their per-line counterparts.
\end{itemize}
\end{statebox}

% ============================================================
% SECTIONS 2–5 — TO BE WRITTEN ONE AT A TIME
% ============================================================

